% answers_short.tex —— 1958 年·华东卷（上海、福建） · 简短答案速查
% 由 tools/make_short_answers.py 自 answers.tex 的 \gksolhead/\gkanswer 提取，
% 并与 1958/1958_华东(上海,福建)_数学真题_答案要点.md 逐题比对后生成，勿手改（改源 MD 后重跑）。

\begin{gknotice}
\textbf{使用说明}：本节只给最终答案便于快速核对，完整推导见解析版。答案为编写者独立推导，\textbf{非官方标准答案}。
\end{gknotice}

\gkpart{题目}
\begin{itemize}[leftmargin=2.6em]
  \item[\textbf{1.}] $m>2\sqrt5$ 或 $m<-2\sqrt5$
  \item[\textbf{2.}] $(x,y)=(1,9)$ 或 $(x,y)=(4,1)$
  \item[\textbf{3.}] 约 $47.42\%$
  \item[\textbf{4.}] (1) $y=2x^2-4x+1$；(2) 无极大值；有极小值，在 $x=1$ 处取得，极小值为 $-1$；(3) $x=1+\dfrac{\sqrt2}{2}$ 或 $x=1-\dfrac{\sqrt2}{2}$；(4) 开口向上的抛物线，顶点 $(1,-1)$，过 $(0,1)$，与 $x$ 轴交于 $\left(1\pm\dfrac{\sqrt2}{2},0\right)$，画法见解析
  \item[\textbf{5.}] 在近岸取基线 $CD=a$，测得 $\angle DCA=\alpha,\angle DCB=\beta,\angle CDA=\gamma,\angle CDB=\delta$，则 $AB=\sqrt{AC^2+BC^2-2\,AC\cdot BC\cos(\alpha-\beta)}$，其中 $AC=\dfrac{a\sin\gamma}{\sin(\alpha+\gamma)}$，$BC=\dfrac{a\sin\delta}{\sin(\beta+\delta)}$
  \item[\textbf{6.}] 见解析版
  \item[\textbf{7.}] $41$ 分米
\end{itemize}
